Spectroscopic Evidence for a 3d10 Ground State Electronic
Configuration and Ligand Field Inversion in [Cu(CF3)4]1−
Richard C. Walroth,† James T. Lukens,† Samantha N. MacMillan,† Kenneth D. Finkelstein,‡
and Kyle M. Lancaster*,†
†Department of Chemistry and Chemical Biology, Baker Laboratory, Cornell University, Ithaca, New York 14853, United States
‡Cornell High Energy Synchrotron Source, Wilson Laboratory, Cornell University, Ithaca, New York 14853, United States
*
S Supporting Information
ABSTRACT: The contested electronic structure of [Cu-
(CF3)4]1− is investigated with UV/visible/near IR spectrosco-
py, Cu K-edge X-ray absorption spectroscopy, and 1s2p
resonant inelastic X-ray scattering. These data, supported by
density functional theory, multiplet theory, and multireference
calculations, support a ground state electronic configuration in
which the lowest unoccupied orbital is of predominantly
trifluoromethyl character. The consensus 3d10 configuration
features an inverted ligand field in which all five metal-localized
molecular orbitals are located at lower energy relative to the
trifluoromethyl-centered σ orbitals.
■ INTRODUCTION
Trifluoromethyl (CF3) substituents profoundly influence the
properties of organic molecules1 and transition metal
complexes.2 Medicinal chemists have recognized the value of
the CF3 group for adding advantageous drug properties
including increased physiological longevity, facile blood−brain
barrier penetration, and enhanced protein−substrate binding
affinities.1 Consequently, the development of chemical
reactions that install CF3 groups on organic frameworks is a
major research goal.3 The desire for late-stage, selective
functionalization has fostered the rapid growth of transition-
metal-mediated trifluoromethylation reactions and has in-
creased interest in preparing discrete, fluorinated organo-
metallic complexes.3c,4 For example, Hartwig and co-workers5
reported the synthesis of Cu(phen)(CF3) (1; phen = 1,10-
phenanthroline; Figure 1). This stable compound is readily
synthesized from inexpensive starting materials and trifluor-
omethylates a wide range of iodoarenes with high functional
group tolerance under mild reaction conditions.
In a recent report, Grushin6 and co-workers described an
improved synthesis of the highly unusual, formally CuIII
complex ion [Cu(CF3)4]1− (2; Figure 1), which is traditionally
prepared with the toxic and explosive Cd(CF3)2 as a
transmetalation reagent. The authors showed that salts of 2
can be easily prepared from CuCl and TMSCF3/KF (TMS =
trimethylsilyl) in air, in a high-yielding, one-step reaction.
These species are air-stable, and in all cases the complex ion
adopts a D2d, flattened tetrahedral geometry. To date, there
have been no reports of 2 exhibiting standard nucleophilic
trifluoromethylation reactivity, although the dearth of studies
may be due in part to the former lack of synthetically accessible
methods for the preparation of 2. Detailed mechanistic studies
by Grushin and co-workers indicate that 2 is an off-path,
catalytically incompetent species in the Cu-mediated oxidative
trifluoromethylation of aryl boronic acids.7
The disputed8 electronic structure of 2 postulated by Snyder9
suggests that 2 should be considered a source of electrophilic,
not nucleophilic, CF3. Although 2 is formally considered a d8
CuIII species, Snyder proposed that it is better described as
physically d10 CuI based on a computational study that gave a
partial charge on Cu of 0.71 and d-orbital populations of 9.7 or
9.4, depending on the level of theory used. With a CuI
configuration, charge balance necessitates that 2 be coordinated
by three CF3
− ligands and one CF3
+ ligand. Owing to
resonance, in a valence bond picture this would best be
described as delocalized CF3
+ character over the molecule, or,
after the manner of Jørgensen,10 CuI surrounded by a (CF3)4
2−
coordination sphere. Snyder’s argument thus describes 2 as
featuring an “inverted” ligand field11 in which the lowest
Received:
October 23, 2015
Published: February 4, 2016
Figure 1. Cu(phen)(CF3) (1; phen =1,10-phenanthroline) is a
versatile trifluoromethylating reagent. To date, trifluoromethylation by
[Cu(CF3)4]1− (2) has not been demonstrated.
Article
pubs.acs.org/JACS
© 2016 American Chemical Society
1922
DOI: 10.1021/jacs.5b10819
J. Am. Chem. Soc. 2016, 138, 1922−1931
Downloaded via FREIE UNIV BERLIN on November 28, 2025 at 14:03:12 (UTC).
See https://pubs.acs.org/sharingguidelines for options on how to legitimately share published articles.


unoccupied molecular orbital (LUMO) exhibits dominantly
ligand orbital admixture, contrary to the classical ligand field in
which the frontier orbitals are dominated by metal d-orbital
admixture (Figure 2). This point was later strongly supported
by Mealli.12 Although such an inversion has been noted in
electronic structure calculations of the LaFe4P12 skutterudite,13
and hinted at in several additional studies,14 an experimentally
verified molecular example remains elusive.
Assignment of oxidation states based on density functional
theory (DFT)-derived partial charges remains controversial,8
and theoretical studies have found that partial charges on metal
centers do not vary with overall charge.11b Moreover, this
particular compound may elude conventional description
because of the counterintuitive electronic and structural effects
imparted by CF3.3c Although theorists have studied 2 for
several years,8,9,11b,15 few experimental electronic structural
characterizations have been offered in support of either ground
state configuration. Herein we report spectroscopic data and
calculations that support d10 as the ground state electronic
configuration of 2.
■ EXPERIMENTAL SECTION
General Considerations. Crystalline samples of 2-PPN [PPN =
bis(triphenylphosphine)iminium] and 2-PPh4 [PPh4 = tetraphenyl-
phosphonium] were provided by V. Grushin, and were prepared
according to a published procedure.16 UV/visible/near IR (UV/vis/
NIR) spectra of solutions of 2 salts dissolved in dichloromethane
(dried using a JC Meyer solvent system) were measured on a Cary
5000 spectrophotometer. CuF2 was purchased from Sigma-Aldrich and
used without further purification.
X-ray Spectroscopy. Samples for X-ray spectroscopy were
prepared as solid mixtures in boron nitride. All mixtures contained
5% by mass Cu. The mixtures were finely ground in an agate mortar,
pressed into 1 mm aluminum spacers, and sealed with 38 μm Kapton
tape.
Cu K-edge X-ray absorption spectroscopy (XAS) spectra of 2-PPN
and 2-PPh4 were measured in transmission mode with N2-filled
ionization chambers at beamline 7−3 of the Stanford Synchrotron
Radiation Lightsource (SSRL) under ring conditions of 3 GeV and
500 mA. Incident X-ray radiation was monochromated by using a
double Si(220) crystal monochromator. Incident energy was calibrated
by using Cu foil as an internal calibrant, with the energy of the first
inflection point set to 8980.3 eV. Samples were maintained at 10 K in
an Oxford liquid He cryostat. Spectra were measured from 8900 to
9300 eV. Four scans were measured and averaged with the
EXAFSPAK17 software package. No spectral changes due to
photodamage were observed after four scans. Background subtraction
was achieved by fitting a polynomial to the pre-edge region and
subtracting this polynomial from the entire spectrum. Spectra were
normalized by fitting a flattened polynomial to the postedge and
normalizing the edge jump at 9000 eV to 1.0.
1s2p resonant inelastic X-ray scattering (RIXS) of 2-PPN and CuF2
were collected at the Cornell High Energy Synchrotron Source
(CHESS) C1 bend magnet beamline under ring conditions of 5.3 GeV
and 100 mA. Incident X-ray radiation was monochromated using a Rh
mirror and sagittal focus double Si(220) monochromator and
calibrated as described above for experiments at SSRL. X-ray emission
was collected, monochromated, and focused with five spherically bent
Si(444) crystals in a Rowland circle geometry by using the CHESS
dual array valence emission spectrometer (DAVES), which has been
described elsewhere.18 X-ray emission was collected with a Pilatus 100
K area detector (Dectris). Two regions of interest (ROIs) were
selected in order to correct X-ray emission intensity (I): the first
(ROI1) contained >95% of the X-ray emission; the second (ROI2)
comprised a region centered on the first ROI but with twice the width
and height. The larger ROI was used to determine the average
background counts per pixel, which was used to estimate the total
background per point on the scan.
=
−
−
−
⎛
⎝
⎜
⎞
⎠
⎟
I
I
I
I
Area
Area
Area
corrected
ROI1
ROI1
ROI2
ROI1
ROI2
ROI1
(1)
X-ray emission was corrected for incident photon flux by dividing
counts from an N2-filled ionization chamber placed just upstream of
the sample.
Emission energy was calibrated by measuring the Kα1 and Kα2 lines
of Cu foil and setting corresponding points on the DAVES Rowland
geometry to 8046.3 and 8026.7 eV, respectively. Photodamage was not
detected in fewer than six scans, so a maximum of six scans were
collected for each spot. Constant 9100 eV incident energy scans of Kα
emission were measured at each new spot on a sample for
normalization. Samples were maintained at ca. 150 K by using an
Oxford N2 cold stream focused at the beam spot. The temperature at
the spot was measured using a thermocouple.
RIXS planes were collected by scanning the emission domain at
constant incident energy. For each compound, 20 incident energies
were chosen along the pre-edge and rising edge regions This strategy is
demonstrated for 2-PPN in Figure S5 and Figure S6. Point densities
were varied to focus collection time on specific regions along the Cu
K-edge. The collection strategies are included as Supporting
Information. Emission intensity was normalized by collecting emission
spectra at a constant 9100 eV incident energy and setting the peak
maximum of the Kα1 peak to 1.
Data were processed with Igor 6.37. Emission spectra were
converted to the energy transfer domain, and individual spectra were
combined to create incident energy (IE), energy transfer (ET), and
intensity tabulated as x, y, and z scatter data. Vertical traces were
generated directly by converting emission spectra to the ET domain.
Horizontal traces were generated by plotting average intensity values
within 0.4 eV of the desired constant energy transfer while ensuring
unique values per incident energy.
Calculations. DFT, time-dependent DFT (TDDFT), DFT-
restricted open-shell configuration interaction with single excitations
(DFT-ROCIS), and spectroscopy-oriented configuration interaction
(SORCI) calculations were performed with version 3.03 of the
ORCA19 software package. All spectra were calculated from crystallo-
graphic coordinates of the 2 complex ion in 2-PPN.16 Calculations did
not employ molecular symmetry. Irreducible representations are
assigned to orbitals by inspection using the symmetry operations in a
D2d character table. Single-point energies were calculated by using
BP86,20 B3LYP,21 M06,22 ωB97,23 ωB97X,23 ωB97X-D3,23 LC-
BLYP,24 and CAM-B3LYP24 functionals. Calculations with hybrid
functionals used the RIJCOSX algorithm to speed the calculation of
Hartree−Fock exchange.25 The CP(PPP)26 basis set was used for Cu
with a special integration accuracy (ORCA Grid7). The scalar
relativistically recontracted def2-TZVP(-f)-ZORA27 basis set with
Figure 2. In a classical Werner-type ligand field, the antibonding
orbital Ψ* is dominated by metal d-orbital character. An inverted
ligand field arises when Ψ* has predominantly ligand orbital character.
Journal of the American Chemical Society
Article
DOI: 10.1021/jacs.5b10819
J. Am. Chem. Soc. 2016, 138, 1922−1931
1923


ORCA Grid4 was used for all other atoms. Calculations included the
zeroth-order regular approximation (ZORA) for relativistic effects28 as
implemented by van Wüllen.29 Solvation was modeled with
COSMO30 in an infinite dielectric.
Electronic absorption and Cu K-edge XAS spectra were calculated
by using TDDFT.31 Spectral predictions used the ωB97X functional.
The CP(PPP) basis set was used for Cu with an integration accuracy
of 7. The def2-TZVP-ZORA basis set was used for all other atoms.
Solvation was again modeled by using COSMO in an infinite
dielectric.
DFT-ROCIS is described elsewhere32 and was used to calculate Cu
L2,3-edge XAS spectra with the ωB97X-calculated electronic structure.
Transitions were restricted to those involving the three Cu 2p orbitals
as donors. Unoccupied orbitals involved in the most intense Cu K-pre-
edge absorptions were chosen as ROCIS acceptor orbitals. CuIII L2,3-
edge XAS was simulated independently with the TT-Multiplets code33
as implemented in the CTM4XAS graphical user interface.34
Multireference character in the ground state of 2 was assessed with
SORCI35 calculations. SORCI was performed on a complete active
space (CAS) for 2 comprising 16 electrons and 10 orbitals
[CAS(16,10)]. The CP(PPP) basis set was used on Cu, and def2-
TZVP(-f)-ZORA was used on all other atoms. The ZORA relativistic
correction was used in all SORCI calculations. As described
elsewhere,35 individual selection was used to ease the computational
burden. The size of the first-order interacting space was reduced with a
threshold: Tsel = 10−6 Eh. A further approximation involved reducing
the reference space through another selection: all initial references that
contributed less than a second threshold (Tpre = 10−5) to the zeroth-
order states were rejected from the reference space. Starting orbitals
were taken from restricted Kohn−Sham orbitals generated via the
aforementioned ωB97X/def2-TZVP(-f)-ZORA calculation.
■ RESULTS AND DISCUSSION
Ligand Field Analysis. The D2d geometry of 2 has been
used to argue against the d10 configuration,8 but regardless of its
configuration, 2 is a 16-electron complex. A σ-only angular
overlap analysis36 comparing inverted D4h and Td ligand field
diagrams reveals that D4h remains the favored geometry despite
the d10 configuration (Figure 3). Additionally, in the d8 limit, 2
should exhibit a moderately absorptive d-d band (ε ca. 100 M−1
cm−1) corresponding to 1A → 1E excitation. No absorption
bands are observed in vis/NIR spectra of solutions of 2-PPN
between 30 000 and 6250 cm−1 (Figure S1). Intense bands due
to charge-transfer absorptions (ε > 20 000 M−1 cm−1) are
found in the UV spectrum to energies greater than 30 000 cm−1
(Figure 4). The burying of the 1A → 1E d-d band beneath the
charge-transfer features implies an unrealistically high value for
crystal field splitting in a four-coordinate Cu complex.
Cu K-edge XAS. Cu K-edge XAS data for 2-PPN were
obtained to complement the UV/vis/NIR data. Cu K-edge
XAS data were also obtained for 2-PPh4the identity of the
counterion has no effect on peak energies (Figure S3). K-edge
XAS has become a standard means of measuring physical
oxidation states in first-row transition metal compounds.37 For
example, Solomon and co-workers38 comprehensively demon-
strated that the pre-edge features in the K-edge XAS spectra of
formally CuIII compounds appear at approximately 8981 eV and
thus are ca. 2 eV to higher energy relative to CuII compounds,
for which pre-edges are typically observed near 8979 eV. This
shift to higher energy exhibited by CuIII complexes is attributed
to an increase in effective nuclear charge commensurate with
higher oxidation state. The increase in the effective charge
lowers the core electrons to a deeper binding energy, which in
turn gives rise to higher-energy pre-edges and rising edges.
The Cu K-edge XAS spectrum for 2-PPN is feature-rich
(Figure 5). Multiple features are resolved from the rising edge,
the inflection point of which occurs at 8997.2 eV. We assigned
the most intense absorption at 8985.6 as the dipole-allowed Cu
1s → 4p transition. This feature is isoenergetic with the 1s →
4p transition in D4h, CuII [CuCl4](nmph)2 (3-nmph; nmph =
Figure 3. σ-Only molecular orbital (MO) diagrams generated via the
angular overlap method for 16-electron complexes with inverted ligand
fields at the Td and D4h geometry limits. Total MO stabilization energy
(MOSE) determinations show that despite the d10 configuration, D4h
geometry is favored in an inverted ligand field picture.
Figure 4. UV/visible spectrum of 2-PPN [PPN = bis-
(triphenylphosphine)iminium] recorded in dichloromethane.
Figure 5. Transmission-mode Cu K-edge X-ray absorption spectros-
copy (XAS) spectrum of solid 2-PPN.
Journal of the American Chemical Society
Article
DOI: 10.1021/jacs.5b10819
J. Am. Chem. Soc. 2016, 138, 1922−1931
1924


methyl(2-phenylethyl)ammonium) and is 1.7 eV lower in
energy compared to analogous transitions in the CuIII species
discussed by Solomon and co-workers.38 The 8985.6 eV feature
is also within the energy range reported for four-coordinate CuI
complexes.39 A considerably weaker feature is found to lower
energy at 8982.0 eV. By convention, the observation of a
feature at this energy supports a CuIII configuration; in these
cases the transition represents a 1s23d8 → 1s13d9 excitation.
However, recent reinterpretation of Cu K-edge XAS data by
Wieghardt and co-workers40 has challenged the diagnostic use
of pre-edge energy as a metric of oxidation state in Cu systems.
DFT as well as ab initio multireference calculations for a
number of formally CuIII complexes have revealed that the d8
electronic configuration makes insignificant contributions to
the ground states of these molecules. These authors40 suggest
that in these cases, the pre-edge excitations are best described
as formally metal-to-ligand charge transfer (MLCT), owing to
the fact that the acceptor orbitals are dominated by ligand
admixture. Contributions from Cu 3d and 4p confer allowed-
ness to these excitations via electronic quadrupole and dipole
contributions, respectively.41 Moreover, weak K-edge XAS pre-
edge features have been observed and assigned in “closed-shell”
d10 CuI systems.42 In these systemssuch as [Cu(bpy)2]+ (bpy
= 2,2′-bipyridine)that bear ligands with unoccupied valence
molecular orbitals (MOs) lying energetically beneath Cu 4p,
such formal metal-to-ligand charge transfer (MLCT) transitions
have been observed near 8981 eV. The energies of these
transitions can be rationally modulated by tuning bpy π-system
energetics, and the transitions are well modeled using TDDFT.
It should be noted that, as distinct from an optical MLCT, the
transition moments giving rise to such features are small, and
the transition densities are effectively metal-localized due to the
highly contracted nature of the donor metal 1s orbital. Such
transitions mirror the pre-edge excitations in ligand K-edge
XAS, where although the transitions may be formally described
as ligand-to-metal charge transfer (LMCT), the contracted
ligand 1s orbitals give rise to small transition moments where
the transition densities are localized to the ligand atoms.43 For
clarity, we will distinguish such charge transfer features as XAS-
MLCT/XAS-LMCT, as distinct from optical MLCT and
LMCT.
In the case of 2, Snyder’s9 postulated CuI configuration
necessitates a “CF3
+” LUMO, which could give rise to a
quadrupole-allowed pre-edge XAS-MLCT feature due to a
1s23d10L → 1s13d10L• excitation in which L represents a CF3 σ-
based MO. Alternatively, the 8982 eV feature could be assigned
to 1s23d104p0 → 1s13d104p1. Excitation to a 1s13d104s1
configuration is forbidden. Both allowed possibilities are
considered below.
1s2p RIXS. The photophysics of Cu K-edge XAS can be
exploited to assign the 8982 eV feature exhibited by 2-PPN.
Excited states with 1s1 configurations rapidly (ca. 1015 s−1)
deactivate through X-ray emission (Figure 6). The deactivation
pathways of highest probability involve Δl = ± 1, 2p → 1s
transitions called Kα emissions. In the d8 CuIII limit (Figure
6a), the final electronic configuration is 1s22p53d9. This
configuration will give rise to multiple nondegenerate electronic
states. Two states, analogous to P1/2 and P3/2 terms in the free
ion limit, arise with a ca. 20 eV splitting owing to Cu 2p spin−
orbit coupling (SOC).44 These states exhibit additional splitting
due to 2p-3d electron−electron exchange. In the d10 limit
(Figure 6b,c), Kα emission from the 1s13d10L• or 1s13d104p1
intermediate states generate either 1s22p53d10L• or
1s23d102p54p1 configurations, the states of which are split by
Cu 2p SOC. In the latter case (Figure 6c), ligand field splitting
of Cu 4p gives two possible excitation/deactivation paths: a
path involving excitation to Cu 4pz, and, to higher energy,
excitation to Cu 4px,y. The exchange interaction between the
Cu 2p and either CF3-localized or 4p-localized electrons is
expected to be vanishingly small owing to the large separation
between radial probability maxima.
These final states are the same as those accessed with Cu L2,3
(2p1/2/2p3/2 → valence/continuum) XAS.45 However, Cu L2,3-
edge XAS is an experimentally demanding soft X-ray technique
that requires ultrahigh vacuum conditions and in which samples
are rapidly photodamaged.46 Moreover, ambiguities can arise in
the assignment of side bands, which may be due to either
electronic structure or X-ray scattering events including
extended X-ray absorption fine structure. Electronic processes
in Cu L2,3 XAS can alternatively be probed with 1s2p RIXS,
which is considered an effective facsimile of L-edge XAS.44,46b,47
RIXS is a hard X-ray, three-dimensional spectroscopy that
presents X-ray emission intensity as a function of incident X-ray
photon energy (IE) and emitted photon energy (EE). The EE
axis is converted by convention into energy transfer ET (ET =
IE − EE). 1s2p RIXS displays energetic separations owing to 2p
SOC, and side bands due to electronic exchange will also be
resolved if the incident energy bandwidth is sufficiently narrow.
Different “slices” of the resulting plane provide comple-
mentary information about the electronic structure of the
compound under study. Plotted diagonal slices give high-
energy-resolution fluorescence-detected (HERFD) XAS spec-
tra. HERFD-XAS reduces uncertainty broadening, thereby
producing XAS spectra with considerably narrower line
widths.48 Plotting slices at constant IE shows the X-ray
emission transitions that deactivate intermediate states
populated by 1s excitations. For example, ET features with
energies that match the L2,3-edge (2p → 3d) XAS maxima are
seen in the 1s2p RIXS analysis of CuF2 at IE resonant with Cu
1s → 3d excitation at 8978.0 eV (Figure S2).45
Figure 6. Qualitative energy diagrams depicting intermediate and final
electronic states accessed by the lowest-energy Cu K-pre-edge
excitations in d8 (a) and d10 (b,c) formulations for 2. The incident
energy (IE) causes excitation from Cu 1s, giving rise to Cu K-edge
XAS features. Subsequent X-ray emission at emission energy (EE)
gives final configurations that match those accessed with Cu L2,3-edge
XAS. The difference (IE − EE) gives energy transfer (ET), the values
of which will match L2,3-edge XAS excitations. In both limits, the final
2p5 configurations will split into two states due to ca. 20 eV Cu 2p
spin−orbit coupling (SOC). In the d8 limit, 2p-3d exchange will cause
additional splitting to give four final states. In the d10 limit where the
initial excitation is to a ligand-centered d-orbital (b), only two final
states will be apparent. In case (c), where the lowest excitation
populates p-orbitals split by the D2d ligand field, four RIXS features
will appear.
Journal of the American Chemical Society
Article
DOI: 10.1021/jacs.5b10819
J. Am. Chem. Soc. 2016, 138, 1922−1931
1925


The 1s2p RIXS plane of 2-PPN obtained with DAVES at
CHESS is shown in Figure 7 (bottom left). The majority of the
X-ray emission intensity is found along parallel diagonals, which
represent final 2p5 states split by Cu 2p SOC amounting to 20.0
eV. Plotting RIXS intensity along these diagonals gives the
HERFD-XAS spectra of 2-PPN (Figure 7, top right), in which
pre-edge and rising-edge features are resolved with full-widths
at half-maximum (fwhm) of 1.5 eV, which is 2-fold narrower
than the width in the corresponding transmission-mode
spectrum (Figure S4). Emission at IEs resonant with the
three observed pre-edge XAS features are plotted in the lower
right of Figure 7. Excitation at 8982.0 eV (red vertical slice)
gives rise to an emission spectrum (Figure 7, bottom right)
with energies determined via the fitting of pseudo-Voigt peaks.
On-diagonal features appear at 934.2 and 954.2 eV. Off-
diagonal bands appear at 937.4 and 958.0 eV. The fwhm for
these features are ca. 2 eV.
Cramer and co-workers44 have shown that 2p-3d and 3d-3d
exchange give rise to well-resolved, ca. 5 eV off-diagonal
splitting in high-spin NiII 1s2p RIXS; low-spin NiII lacks 3d-3d
exchange splitting. However, neither interaction is the origin of
the side bands in the 1s2p RIXS spectrum of 2-PPN. Another
origin would be 3d-3d exchange. However, salts of 2 are
diamagnetic,16 necessitating either a closed-shell or singlet
diradical ground state configuration. In the case of either a
classical or inverted ligand field, the highest MO with Cu d-
character is nondegenerate. In the diradical case, 3d-3d
exchange would energetically favor a triplet ground state over
a singlet. Moreover, the ground state would have orbital
degeneracy and be subject to a Jahn−Teller distortion. While
the latter consequence need not necessarily manifest in
observable geometric distortion, the diamagnetism of 2 rules
out open-shell configurations. Consequently, no side bands will
arise due to 3d-3d exchange in the RIXS of 2-PPN. While
additional features are observed at 937.4 and 958.0 eV, it
should be noted that these bands are isoenergetic with the
emission features arising from 8985.6 eV IE (blue vertical slice).
Plotting constant ET at 937.4 and 958.0 eV (denoted by the
dashed lines in Figure 7) gives peak maxima at 8985.6 eV IE
(see Figure 7, top left blue spectrum). Thus, these features do
not originate from excitation at 8982.0 eV, but rather arise from
the decay of the higher-energy intermediate state with
configuration 1s14p1 that is adventitiously accessed because
IE bandwidth exceeds that of the 8982 eV pre-edge feature.
This causes excitation into the state corresponding to the next
XAS feature at 8985.6 eV. Plotting emission intensity at ET
values that would be consistent with 2p-3d exchange (vide
infra) as a function of IE shows that all emission intensity arises
due to 8985.6 eV IE and that none of this intensity originates
from 8982 eV IE. The lack of residual off-diagonal intensity
arising from 8982 eV IE is consistent with vanishingly small Cu
2p-3d exchange, although distinction between excitation to
1s13d10L• versus 1s13d104p1 remains at this point ambiguous.
Regardless, the 1s2p RIXS spectrum for 2-PPN is consistent
with a ground state electronic configuration approaching the
d10 limit.
We assign the 934.2 and 954.2 eV features as the L3-and L2-
edge maxima, respectively, for 2-PPN. The energies of these
peaks are both ca. 2 eV higher than the corresponding values
found during L-edge XAS studies of CuIII species by Solomon
and co-workers.49 Together with the high energy found for the
first pre-edge band in the Cu K-edge XAS spectrum, this shift
suggests a LUMO with an energy placed above the Cu 3d
manifold, also consistent with the inverted ligand field
framework.
DFT Calculations. Single-point DFT calculations were
carried out with ORCA19 to produce an MO description of the
ground state electronic configuration of 2. Solomon and co-
workers have noted that calculations of Cu covalency are highly
functional-dependent. These authors benchmarked several
functionals in calculations of the electronic structure of the
thoroughly characterized D4h CuII complex ion CuCl4
2− (3), for
which a self-consistent, experimental value of 0.62 ± 0.02
electrons has been determined for the spin population of the
singly occupied MO (Cu 3dx
2−y
2).50 These authors observed
that BP86 returns an inverted ligand field for 3, which is an
electronic structure that is inconsistent with spectroscopic
characterization.
We carried out single-point calculations of 3 with various
density functionals while keeping constant the use of the
CP(PPP) basis set on Cu and the segmented all-electron
relativistically recontracted def2-TZVP(-f)-ZORA basis set on
Cl (Table 1). Our test set of functionals included several
modern range-separated functionals.23 Löwdin spin population
analyses were used to estimate Cu−Cl covalency.51 Whereas
BP86 returns an inverted ligand field for 3, B3LYP, M06 and
the range-separated functionals produce a bonding scheme that
is more consistent with experiment. We noted that ωB97X and
its dispersion-corrected ωB97X-D3 variant gave splendid
agreement with the experimental value, yielding values of
0.61 and 0.60 electrons, respectively. Cognizant of the
functional sensitivity of Cu covalency, we employed BP86,
Figure 7. 1s2p resonant inelastic X-ray scattering (RIXS) plane of
solid 2-PPN is plotted in the lower left. Constant ET slices with colors
corresponding to solid horizontal lines on the RIXS plane are plotted
in the upper left. Constant IE slices with colors corresponding to the
vertical solid lines on the RIXS plane are plotted in the lower right.
Points on the Cu K-edge XAS spectrum corresponding to the IE slices
are marked with color-coded asterisks. A diagonal slice corresponding
to the black solid line gives the HERFD XAS spectrum of 2-PPN,
which is plotted in the upper right.
Journal of the American Chemical Society
Article
DOI: 10.1021/jacs.5b10819
J. Am. Chem. Soc. 2016, 138, 1922−1931
1926


B3LYP, and ωB97X in DFT calculations of the electronic
structure of 2.
Frontier MO diagrams for 2 produced at the BP86/def2-
TZVP(-f)-ZORA, B3LYP/def2-TZVP(-f)-ZORA, and ωB97X/
def2-TZVP(-f)-ZORA levels of theory are shown in Figure 8
and orbital energies and coefficients are given in Table 2. In
accord with Snyder’s9 earlier finding, the b2-symmetry σ*
LUMO (2b2) is dominated by CF3 σ character in all cases.
These results are generated using crystallographic coordinates
for 2 in 2-PPN, and geometry optimizations using either BP86
or B3LYP do not influence the outcome (see Tables S1 and
S2). The Cu 3dxy contributions to 2b2 range from 29% (BP86)
to 33% (ωB97X). CF3 σ MOs contribute 64% (BP86) to 59%
(ωB97X). Residual contributions to 2b2 are from Cu 4pz. The
three highest occupied MOs calculated via B3LYP and ωB97X
are principally of CF3 parentage, and to deeper binding energy
are found the five fully occupied Cu 3d orbitals with an
ordering that is the inverse of classical Td ligand field-splitting.
The BP86 results largely preserve the ligand field inversion,
although the ligand-centered a1 orbital is calculated to lower
energy than the metal-centered a1.
Spectral Simulations via TDDFT and ROCIS. All
employed density functionals produce a consensus electronic
structure in which the conventional ordering of metal-centered
and ligand-centered MOs is inverted. The aforementioned
electronic structure calculations using crystallographic coor-
dinates were used as starting points for spectral predictions.
Results of these calculations are essentially identical, so only
spectral calculations employing ωB97X will be discussed; the
results from the BP86 and B3LYP calculations are included in
the SI.
Table 1. Calibration Calculations of Cu Covalency in 3a
functional
% Cu 3db
BP86
0.47
B3LYP
0.56
M06
0.55
ωB97X
0.61
ωB97X-D3
0.60
CAM-B3LYP
0.59
LC-BLYP
0.54
experimentc
0.62 ± 0.02
aCalculations used the def2-TZVP(-f)-ZORA basis for Cl and
CP(PPP) for Cu. bSpin population on Cu. cSee ref 49.
Figure 8. Frontier MO diagrams of 2 at the BP86/def2-TZVP(-f)-ZORA, B3LYP/def2-TZVP(-f)-ZORA, and ωB97X/def2-TZVP(-f)-ZORA levels
of theory. The CP(PPP) basis set was used on Cu in all calculations. Cu-localized MOs are shown in black, and ligand-localized orbitals are gray. The
b2 bonding and antibonding MOs from the ωB97X calculation are plotted at an isolevel of 0.03 au.
Table 2. Molecular Orbital Energies and Coefficientsa for 2b
BP86
orbital
energy
(eV)
% Cu
3d
% Cu
4s
% Cu
4p
% CF3 σ
2b2 (LUMO)
−3.3
29.8
0.0
4.5 (z)
63.8
2e (2-fold degenerate)
−6.6
4.8
0.0
6.4 (x)
76.8
6.4 (y)
2a1
−7.4
4.8
10.5
0.0
26.7
1e (2-fold degenerate)
−8.2
59.5
0.2 (x)
9.3
0.2 (y)
1b1
−8.5
88.7
0.0
0.0
8.9
1a1
−8.9
88.7
3.6
0.0
53.5
1b2
−9.4
89.6
0.0
0.0
37.9
B3LYP
orbital
energy
(eV)
% Cu
3d
% Cu
4s
% Cu
4p
% CF3 σ
2b2 (LUMO)
−2.8
32.0
0.0
4.5 (z)
61.5
2e (2-fold degenerate)
−7.7
1.9
0.0
6.7 (x)
78.7
6.7 (y)
2a1
−9.0
35.9
12.7
0.0
46.4
1e (2-fold degenerate)
−10.3
90.2
0.0
0.0
8.5
1a1
−10.7
62.0
1.1
0.0
31.6
1b1
−10.7
82.8
0.1
0.0
15.9
1b2
−10.8
56.6
0.0
0.0
38.6
ωB97X
orbital
energy
(eV)
% Cu
3d
% Cu
4s
% Cu
4p
% CF3 σ
2b2 (LUMO)
−0.7
34.1
0.0
4.7 (z)
58.9
2e (2-fold degenerate)
−10.4
1.2
0.0
6.9 (x)
78.8
6.9 (y)
2a1
−12.1
27.9
12.6
0.0
53.3
1e (2-fold degenerate)
−13.5
86.9
0.0
0.0
12.4
0.0
1b2
−13.8
52.2
0.0
0.0
42.6
1a1
−13.9
66.4
0.5
0.0
28.5
1b1
−13.9
80.1
0.0
0.0
19.1
aSummed Löwdin orbital coefficients. bCalculations on crystallo-
graphic coordinates using the CP(PPP) basis set on Cu and def2-
TZVP(-f)-ZORA on all other atoms.
Journal of the American Chemical Society
Article
DOI: 10.1021/jacs.5b10819
J. Am. Chem. Soc. 2016, 138, 1922−1931
1927


The lowest-lying spin-allowed excited state calculated via
TDDFT was found at 37 089.4 cm−1, which agreed with the
UV/vis/NIR results, and it arises due to charge-transfer
excitations (Figures S10). The dominant contributing ex-
citations to the ca. 35 000 cm−1 absorption feature are 2-fold
degenerate, 1A1 → 1E MLCT transitions. No d-d bands were
predicted. TDDFT calculations of the Cu K-edge XAS results
accurately reproduced the relative energy separations among
the first three pre-edge bands resolved in the HERFD data with
a correlation coefficient of R2 = 1.0 (Figure 9a,b). The 8982 eV
pre-edge feature is predicted by TDDFT calculations to be Cu
1s → 2b2 σ*(CF3). Under D2d symmetry, 4pz will mix with this
MO, in accord with the relatively high intensity of this pre-edge
transition. The feature at 8985.6 eV arises due to Cu 1s → 4pz
excitations. The feature at 8989.5 eV comprises numerous
excitations. Major contributions arise due to excitations from
Cu 1s to CF3 π* orbitals with small (ca. 5%) Cu 3d admixture.
The TDDFT-calculated Cu K-edge spectra were used to
guide DFT-ROCIS32 calculations of Cu L2,3 XAS findings to
model the 1s2p RIXS data. For the DFT-ROCIS calculations,
the Cu 2p orbitals were chosen as donors, and three valence
MOs participating in the 8982.0, 8985.6, and 8989.5 eV
excitations were chosen as acceptor orbitals. The calculated
energies of the three pairs of Cu 2p SOC-split features
correlated with R2 = 1.0 to experiment (Figure 9c,d).
Importantly, no side-bands due to 2p-3d exchange were
observed in the calculated RIXS features arising due to the
8982.0 eV IE excitation.
Relative L2,3 peak areas are well-reproduced by the DFT-
ROCIS procedure for all functionals investigated.52 The ratios
of peak areas in L2,3-edge XAS are conventionally expressed in
terms of the branching ratio, which is given as I(L3)/[I(L3) +
I(L2)]. Thole and Van der Laan53 have shown that this
branching ratio is linearly related to valence SOC. In the
present case, this is the magnitude of Cu 3d SOC and, thus, the
branching ratio reflects the relative contributions of Cu 3d and
CF3 σ* to the orbital probed by emission with 8982 eV
resonant excitation. The branching ratio resulting from DFT-
ROCIS calculations is 0.64 for all functionals, in good
agreement with the experimental value of 0.69. This branching
ratio is consistent with an acceptor orbital involved in the 8982
eV excitation for which Cu 3d makes a minor contribution
(vide infra).
Intensities for RIXS ET bands corresponding to 8985.6 and
8989.5 eV IE were severely underestimated with this method.
However, this outcome is readily explained: these features arise
in the RIXS slices due to 2p → 1s demotions and are formally
|Δl| = 1, thus obeying atomic selection rules. The DFT-ROCIS
method calculates these bands as L2,3 excitations. Because the
8985 eV band is Cu 2p → Cu 4p with a Δl of 0, the
corresponding RIXS features are formally forbidden and thus of
low intensity. The large disparity in intensities between the
RIXS features arising from 8982 and 8985.6 eV excitation
supports assignment of the 8982 eV excitation as XAS-MLCT
from Cu 1s to 2b2 σ*(CF3). The band at 8989.5 eV contains
contributions from CF3 π* orbitals with ca. 5% 3d
contributions. The corresponding L2,3 excitation will be weak
due to this small d-contribution.
X-ray Spectral Simulations via Multiplet Theory.
Multiplet calculations33 were carried out to evaluate the
magnitude of Cu 2p-3d exchange splitting that would be
apparent when accessing 2p53d9 states with 1s2p RIXS in the
Figure 9. (a) High-energy-resolution fluorescence-detected (HERFD) XAS spectrum of 2-PPN (black) overlaid on the energy-corrected time-
dependent density functional theory calculated spectrum (red). (b) Calculated transition energies were empirically adjusted according to the strong
correlation to experimental peak energies. (c) Calculated, energy-corrected 1s2p RIXS ET features of 2-PPN arising at IEs corresponding to the
three lowest energy pre-edge XAS features. (d) Calculated RIXS ET peak energies correlate strongly with experimental values. All spectral
calculations were initiated from the ωB97X/def2-TZVP(-f)-ZORA solution.
Journal of the American Chemical Society
Article
DOI: 10.1021/jacs.5b10819
J. Am. Chem. Soc. 2016, 138, 1922−1931
1928


CuIII d8 limit. An Oh crystal field was chosen with Δo estimated
at 0.5 eV from DFT. The value of Fdd was set to 0 to reflect
diminished 3d-3d exchange due to the closed-shell ground state
configuration of 2. Values of 2p-3d exchange were scaled from
100% (d8 limit) to 0% (d10 limit) by concomitantly varying
reduction factors of the Slater Fpd and Gpd integrals from 1.0 to
0.0. A pseudo-Voigt broadening was chosen with a total fwhm
value of 2.0 eV comprising equal Gaussian and Lorentzian
contributions. This value was chosen based on experimental
line widths. The results of these calculations are plotted in
Figure 10.
The multiplet calculations show that 2p-3d exchange
manifests in the d8 limit as a 2.4 eV split of the L3 band into
two features of approximately equal intensity. At reduction
factors of 40% and lower (approaching d10), these features
coalesce to a single band. The calculation with 40% reduction
of the 2p-3d integrals gives an L3 peak at 934.2 eV and an L2
peak at 954.5 eV in excellent agreement with experimental
values. However, the relative peak areas do not agree with
experiment for any value of integral reduction.
A second set of multiplet calculations was performed, this
time holding Fdd at 0.0, Fpd at 0.4, and Gpd at 0.4. Valence SOC
(the 3d SOC) was scaled from 0% to 100%, and the resulting
spectra are plotted in Figure 11. Shown in the inset of Figure 11
is the linear relationship between the branching ratio and the
percent valence SOC. Interpolation of the experimental value
of the branching ratio (0.69) found for 2 shows that the data
are best modeled with the valence SOC reduced to 15%, in
further accord with a largely CF3-centered LUMO.
SORCI Calculations. SORCI calculations were carried out
to evaluate multireference character in the ground state of 2.
SORCI calculations were initiated from a CAS containing the
16 valence electrons in σ orbitals. The eight orbitals harboring
these electrons, as well as the LUMO and LUMO+1, were
chosen for the CAS for a final CAS(16,10) reference. SORCI
calculations were initiated with orbitals from the ωB97X/def2-
TZVP(-f)-ZORA DFT calculation. Ten singlet, triplet, and
quintet states were calculated, and the chosen CAS captured
>90% of the reference configurations. Both variations of the
SORCI procedure revealed that the ground state of 2 is
effectively single-reference and dominated by the closed-shell
configuration determined with DFT (Table S7). The reference
weight of this configuration was 0.86. Singlet diradical
configurations amount to less than 5% of the ground state
references. The lowest-lying triplet and quintet states are
calculated at 3.2 and 8.9 eV above the ground state,
respectively. The inverted ligand field MO scheme for 2 was
conserved after the configuration interaction (CI) treatment
(Figure S13). The UV/vis spectrum generated by using the ten
states of each spin multiplicity is provided in Figure S14. This
spectrum is in good agreement with the experimental UV/vis/
NIR data and again features a lack of d-d transitions.
Extension to Additional Cu-CF3 Complexes. The
aforementioned analyses indicate that DFT calculations at the
ωB97X/def2-TZVP(-f)-ZORA level successfully model the
electronic structure of 2. These calculations were extended to
1 as well as to Cu(bpy)(CF3)3 (4) and Cu(edtc)(CF3)2 (5;
edtc = diethyldithiocarbamate; Figure 12). Similar to 2, 4 and 5
have to date been described as formally d8 CuIII species.
However, ligand field inversion resurfaces, with the LUMOs of
both molecules exhibiting electronic vacancy of CF3-centered
MOs. The five Cu-localized MOs are filled, giving a d10
configuration. By contrast, the LUMO of 1 is phen π*. All
CF3-localized MOs are filled, and in light of these comparisons,
the reactivity of 1 as a source of nucleophilic “CF3
−” is readily
apparent, whereas the presence of a formally “CF3
+” species on
2, 4, and 5 suggests that these molecules could serve as unique,
metal-based sources of electrophilic CF3. Such reactivity studies
are under way in our laboratory.
■ SUMMARY AND OUTLOOK
The assembled spectroscopic data and supporting calculations
establish a limiting d10 configuration for 2. This result
necessitates an inversion of the MO ordering from the classical,
Werner-type ligand field picture. Formally, the coordination
sphere of 2 should exhibit electrophilic “CF3
+” character. Our
results showcase the utility of 1s2p RIXS for the assignment of
Figure 10. TT-multiplet simulations of Cu L2,3-edge XAS data from
the d10 to d8 limits in 2p-3d exchange. An octahedral crystal field with a
Δo of 0.5 eV was used for all calculated spectra. 3d-3d exchange was
abolished by setting Fdd to 0.0. 2p-3d exchange was modulated by
varying Fpd and Gpd together from 0.0 (d10 limit) to 1.0 (d8 limit).
Pseudo-Voigt line broadening with 2 eV full-widths at half-maximum
(fwhm) was applied to each spectrum.
Figure 11. TT-multiplet simulations of Cu L2,3-edge XAS data from
the d10 to d8 limits in valence spin−orbit coupling (SOC). An
octahedral crystal field with a Δo of 0.5 eV was used for all calculated
spectra. 3d-3d exchange was abolished by setting Fdd to 0.0. 2p-3d
exchange was held at 40% reduction with Fpd and Gpd values set to 0.4.
The value of valence (3d) SOC was scaled from 0% to 100%. A
pseudo-Voigt line broadening with 2 eV fwhm was applied to each
spectrum. The inset shows the linear correlation between valence SOC
and the L3/L2 branching ratio defined as I(L3)/[I(L3) + I(L2)]. A
vertical line marks the branching ratio of 2, corresponding to 15% 3d
SOC.
Journal of the American Chemical Society
Article
DOI: 10.1021/jacs.5b10819
J. Am. Chem. Soc. 2016, 138, 1922−1931
1929


XAS pre-edge features, and they show that XAS-MLCT features
can arise at energies once considered diagnostic of CuIII. The
ROCIS-DFT method predicts 1s2p RIXS ET features with
relative energy separations that scale to splendid agreement
with experiment. Extending our experimentally validated
calculations to related formally CuIII compounds shows that
the d8 description is, in several cases, inconsistent with
experimentally validated electronic structure calculations.
These results, combined with those of the aforementioned
studies by Wieghardt and co-workers,40 indicate that
reevaluation of the electronic structures of high-valent Cu
speciesparticularly purported CuIV species including
Cs2CuF6
54 and recently reported Cu corrolato monoca-
tions55is merited. Such studies, as well as investigations
into both the physical origin and the reactivity consequences of
the ligand field inversion exhibited by 2 and related compounds
are forthcoming.
■ ASSOCIATED CONTENT
*
S
Supporting Information
The Supporting Information is available free of charge on the
ACS Publications website at DOI: 10.1021/jacs.5b10819.
NIR data for 2-PPN, Cu K-edge XAS data for 2-PPh4,
1s2p RIXS data for CuF2, 1s2p RIXS data collection
strategy, representative ORCA input files, ORCA input
coordinates, DFT orbital plots for 2, DFT orbital
energies and admixture coefficients following geometry
optimization of 2, DFT orbital energies and admixture
coefficients for 1, 4, and 5, TDDFT calculated UV/vis/
NIR spectra of 2, SORCI results, and SORCI-calculated
UV/vis/NIR spectrum of 2. (PDF)
■ AUTHOR INFORMATION
Corresponding Author
*kml236@cornell.edu
Notes
The authors declare no competing financial interest.
■ ACKNOWLEDGMENTS
We thank Vladimir Grushin for supplying the salts of the title
complex ion, and we thank Roald Hoffmann for valuable
discussions. K.M.L. thanks the Cornell University College of
Arts and Sciences for startup funding and support from the
National Science Foundation (NSF) in the form of a CAREER
award (CHE-1454455). This work is based in part on research
conducted at CHESS, which is supported by NSF and the
National Institutes of Health/National Institute of General
Medical Sciences (NIH/NIGMS) under NSF award DMR-
0936374. Transmission-mode XAS was performed at SSRL,
which is supported by the U.S. Department of Energy, Office of
Science, Office of Basic Energy Sciences under Contract No.
DE-AC02-76SF00515. The SSRL Structural Molecular Biology
Program is supported by the Department of Energy’s Office of
Biological and Environmental Research, and by NIH/NIGMS
(including P41GM103393). The content of this publication is
solely the responsibility of the authors and does not necessarily
represent the official views of NIGMS or NIH.



### The intial xyz structure that are used in this paper are in the following 
### unit are in angstrom



Cu(phen)(CF3) (1)
C 1.242750 -2.370150 -1.541270
 C 1.600590 -3.069920 -2.709180
 C 1.303170 -2.525740 -3.946110
 C 0.644570 -1.275930 -4.012340
 C 0.323590 -0.646880 -2.782430
 N 0.620030 -1.188460 -1.568040
 C -0.346920 0.626060 -2.784240
 C -0.682890 1.243920 -4.015810
 C -0.345920 0.583330 -5.244250
 C 0.291840 -0.626970 -5.242580
 C -1.248130 2.362250 -1.547770
 C -1.621060 3.050900 -2.717540
 C -1.340210 2.494610 -3.952970
 H 1.467400 -2.778280 -0.554960
 H 2.107240 -4.030570 -2.621810
 H 1.569790 -3.046880 -4.867150
 H -0.609940 1.068440 -6.185490
 H 0.543620 -1.121070 -6.182500
 H -1.457980 2.780700 -0.562540
 H -2.125870 4.012740 -2.632740
 H -1.618550 3.007120 -4.875390
 N -0.627010 1.179650 -1.571420
 Cu 0.000000 0.000000 0.000000
 C -0.000000 -0.000000 1.908520
 F -0.098790 1.259370 2.542630
 F -1.059410 -0.719790 2.497370
 F 1.135670 -0.554120 2.538000


[Cu(CF3)4] 1– (2)

C 1.98590 -0.02661 -0.12375
F 2.49430 0.95471 -0.95352
F 2.48870 -1.19316 -0.66621
F 2.63413 0.13295 1.07811
C -1.96175 -0.02906 -0.33812
F -2.39764 0.97970 -1.17457
F -2.37896 -1.17826 -0.98171
F -2.73705 0.06843 0.79234
C -0.03479 -1.93631 0.46443
F -1.15991 -2.30629 1.17722
F 1.00336 -2.32228 1.28910
F 0.01347 -2.78502 -0.61588
C -0.00000 1.99004 0.00000
F -1.11201 2.53583 0.61195
F 0.05349 2.55828 -1.25017
F 1.05252 2.54117 0.70512
Cu 0.00000 0.00000 0.00000

[Cu(CF3)4]
1– (2)

C -1.92955 -0.02958 -0.48028
F -2.31320 -1.19061 -1.12397
F -2.28808 0.96228 -1.37266
F -2.78865 0.10306 0.58522
C 1.98957 -0.02827 0.00375
F 2.52459 -1.19590 -0.50496
F 2.54538 0.95202 -0.79671
F 2.56595 0.13352 1.24155
C -0.00000 1.98948 0.00000
F 0.99290 2.53269 0.79240
F -1.15713 2.53701 0.52104
F 0.15832 2.56667 -1.23772
C -0.05845 -1.93004 0.47805
F 0.92105 -2.29767 1.38064
F 0.07664 -2.79034 -0.58630
F -1.22918 -2.30183 1.11008
Cu 0.00000 0.00000 0.00000


[Cu(CF3)4]
1– (2)

C 1.99528 -0.00000 -0.00000
F 2.54002 1.02117 -0.72990
F 2.50329 -1.12814 -0.59053
F 2.59252 0.07334 1.21869
C -1.91916 -0.06180 -0.55767
F -2.14450 0.81330 -1.58964
F -2.31895 -1.26713 -1.06368
F -2.84188 0.25120 0.39128
C -0.03244 -1.91831 0.50475
F -1.14869 -2.22110 1.24132
F 1.00212 -2.30690 1.30862
F -0.02660 -2.78897 -0.54104
C -0.04853 1.98931 -0.00013
F -1.20796 2.50523 0.51385
F 0.10086 2.57627 -1.22057
F 0.92098 2.53335 0.79907
Cu 0.00000 0.00000 0.00000


D4h [CuCl4]

Cu 0.00000 0.00000 0.00000
Cl -0.00000 -2.28445 0.00000
Cl 2.20675 -0.00219 -0.46078
Cl 0.00000 2.28445 -0.00000
Cl -2.20675 0.00219 0.46078




Cu(bpy)(CF3)3 (4)

Cu 0.000000 0.000000 0.000000
 N 2.020690 0.016350 0.017730
 C 2.751970 0.049420 -1.109360
 H 2.199270 0.065890 -2.048490
 C 4.143670 0.063020 -1.088460
 H 4.698350 0.089920 -2.025150
 C 4.790340 0.042070 0.147530
 H 5.878820 0.052500 0.205730
 C 4.029320 0.008080 1.313300
 H 4.521280 -0.007730 2.283870
 C 2.631840 -0.005150 1.226980
 N 0.418280 -0.039040 2.157580
 C 1.744400 -0.040780 2.415670
 C 2.230060 -0.074690 3.729220
 H 3.299700 -0.077930 3.930590
 C 1.321920 -0.105670 4.785830
 H 1.681200 -0.132700 5.814810
 C -0.045510 -0.102130 4.509230
 H -0.788420 -0.125370 5.305660
 C -0.453250 -0.068770 3.175780
 H -1.510990 -0.065900 2.909060
 C 0.000000 1.999160 -0.000000
 C -1.825140 0.006390 -0.756340
 C 0.024570 -1.996610 -0.068730
 F 0.187010 2.524760 -1.257870
 F 1.023660 2.530560 0.759740
 F -1.135360 2.583280 0.501740
 F -2.732560 0.016740 0.279580
 F -2.132250 1.085990 -1.541240
 F -2.151610 -1.073110 -1.532970
 F -1.093520 -2.607250 0.440950
 F 0.188020 -2.481520 -1.345820
 F 1.070440 -2.541440 0.650900


Cu(edtc)(CF3)2 (5)

Cu -0.000000 0.000000 0.000000
 S -2.108050 0.462060 -0.652690
 S -0.000000 2.254890 -0.000000
 F 0.781800 -2.655870 -0.595460
 F -1.373040 -2.358030 -0.708990
 F -0.396440 -2.436600 1.247740
 F 2.159370 -1.123560 1.453970
 F 2.664910 -0.444690 -0.574420
 F 2.415870 1.001340 1.051180
 N -2.446910 3.157970 -0.764010
 C -1.653510 2.131860 -0.514030
 C -1.965020 4.550520 -0.639910
 C -1.467630 5.115430 -1.969160
 C -3.854450 2.963790 -1.172410
 C -4.817740 3.023720 0.011970
 C 1.896550 -0.158400 0.524350
 C -0.247320 -1.958740 -0.028290
 H -1.170060 4.565760 0.118280
 H -2.804060 5.140680 -0.247290
 H -0.613660 4.537860 -2.347720
 H -1.143720 6.154580 -1.819920
 H -2.260200 5.113050 -2.729890
 H -3.925760 1.997840 -1.691090
 H -4.081450 3.747490 -1.907940
 H -5.846490 2.903320 -0.353870
 H -4.754360 3.987440 0.535400
 H -4.609400 2.219400 0.729660